\begin{frame}{Дополнительно: почему Drosophila и почему ECAP не мера боли}
\textbf{Почему Drosophila.} Для взрослой мухи доступны коннектом и модели нейронных цепей, проверенные по активности. Переносятся только начальные параметры временного модуля; анатомическое соответствие видов не предполагается.
\SlideGap
\textbf{Почему ECAP не мера боли.} ECAP отражает рекрутирование волокон в измерительном канале и зависит от условий регистрации. Ноцицептивная реакция мухи не тождественна субъективной боли человека; клинические выводы требуют связанных продольных данных SCS.
\SourceNote{Shiu et al., Nature, 2024. DOI: \href{https://doi.org/10.1038/s41586-024-07763-9}{10.1038/s41586-024-07763-9}. Lappalainen et al., Nature, 2024. DOI: \href{https://doi.org/10.1038/s41586-024-07939-3}{10.1038/s41586-024-07939-3}. Anaya et al., Neuromodulation, 2020. DOI: \href{https://doi.org/10.1111/ner.12965}{10.1111/ner.12965}.}
\end{frame}
